\apendice{Anexo de sostenibilización curricular}

\section{Introducción}
Este capitulo se incluye una reflexión personal sobre los distintos aspectos de la sostenibilidad que se abordan en este proyecto. A continuación expongo mi reflexión sobre los principios de sostenibilidad aplicados durante la realización de este Trabajo Fin de Grado, siguiendo las directrices de la \href{https://www.crue.org/wp-content/uploads/2020/02/Directrices_Sosteniblidad_Crue2012.pdf}{CRUE}.


\section{Competencias de sostenibilidad}

\subsection{Pensamiento enfocado a la sostenibilidad}
Al planificar el proyecto analizamos múltiples factores a considerar y realizamos una cadena completa: dispositivo de captura, equipos con \textit{hardware} limitado,interfaz de usuario, colectivo beneficiario y mantenimiento futuro. Ese mapa de me permitió priorizar cambios que aportaban más valor social sin aumentar la huella ambiental.

\subsection{Responsabilidad social y accesibilidad}
Diseñé la interfaz con contraste alto de colores, navegación por teclado y mensajes cortos sencillos para facilitar a los usuarios su navegación por la aplicación web de reconocimiento ASL. Además, todos los procesos se ejecutan localmente, no se guardan imágenes ni se envían datos a terceros, lo que protege la privacidad de los usuarios. 

\subsection{Accesibilidad del proyecto}
Siempre hemos tenido en cuenta a la hora de realizar nuestro proyecto una aplicación ligera y accesible para todos. Desde un primer momento se quiso resaltar la idea de que dispositivos computacionalmente muy limitados y pequeños actualmente son capaces de correr modelos para realizar clasificación de imágenes con modelos optimizados y cuantizados, Demostrando que la clasificación visual no está limitada únicamente a equipos de alto rendimiento.

\subsection{Uso responsable de recursos y software libre}
Para promover la reutilización, todo el proyecto se ha desarrollado con dependencias \textit{open-source} apoyándonos en librerías con licencias permisivas (MIT, Apache 2.0, BSD). este enfoque evita costes de licencias y facilita el uso de estas en nuestro proyecto. Además, minimizar la instalación de software adicional y reutilizar módulos existentes contribuye a un menor consumo de recursos y favorece un ciclo de desarrollo más ágil, colaborativo y sostenible.




\section{Conclusión}
La elaboración de este proyecto sobre la detección de gestos de lenguaje ASL me ha demostrado que la sostenibilidad no es un extra, sino algo que se puede (y se debe) tener presente desde la primera línea de código.

Al optimizar el modelo para que funcione en una Raspberry Pi y al utilizar librerías de código abierto además de cuidar la accesibilidad de la interfaz, he visto cómo decisiones pequeñas marcan la diferencia: mayor rendimiento en sistemas con bajos recursos, menos barreras para quien quiera usar el sistema y más facilidad para implementar actualizaciones en un futuro.

Sobre esto último, a corto plazo quiero realizar modificaciones presentadas como líneas a futuro en la memoria. Así cierro un proyecto que, además de reconocer lenguaje ASL, demuestra que se puede hacer IA de forma eficiente, abierta y útil para todos.